\scanpage{30}
% archive foliation: 29
qui sont $\varprojlim$ du type envisagé de
foncteurs représentables (dans le
cas resp., on utilise, pour y donner
un sens à cet énoncé --- dans $\widehat A_{ab}$ ---
la structure ``\unsure{$k$|$\mathbf Z$}-additive'' de $A$
héritée par son plongement dans $C$
additive). Lorsque de plus $C$
est stable par ce type de $\varprojlim$
(p.ex. des produits finis), alors
on a $\simeq$ une \uline{équivalence} entre
$C$ et la cat. des foncteurs
$A^\circ\longrightarrow\mathrm{Ens}$ (resp. $A^\circ\longrightarrow\mathrm{Ab}$) qui sont
de telles $\varprojlim$ de foncteurs représen-
tables.
\deleted{[Si p.ex. $A$ est réduite à
un objet, \unsure{on trouve donc $C$}
($C$ additive, en resp., si $\forall$ objet}

Si p.ex. $C$ est additive, donc $A$
\unsure{induit sur $A$} une structure ``\unsure{$k$|$\mathbf Z$}-
additive'', et si $\forall$ objet de $C$
est isomorphe à une somme
directe \inserted{finie} d'objets de $A$, alors
$C\simeq$ sous-cat. pleine de $\widehat A_{ab}$
(foncteurs additifs contravariants
de $A$ dans $\mathrm{Ab}$) qui sont des
foncteurs isom. à une somme
directe finie de foncteurs représ.
